\begin{tablalafrox}{Diccionario del Producto 3: Plataforma base. Fuente: según cita por cuenta en cada encabezado.}{tab:T14-dic-prod-3}
  {L{1.2cm} Y Y Y C{0.9cm} L{2.3cm}}
  {\cab{Código} & \cab{Paquete de trabajo} & \cab{Entregable} & \cab{Criterio de aceptación} & \cab{Resp.} & \cab{Período}}
  \multicolumn{6}{>{\columncolor{lafroxBg}\RaggedRight\arraybackslash}p{\dimexpr\textwidth-2\tabcolsep\relax}}{\textbf{Cuenta 3.1 — Ambientes y entrega continua.} Esta cuenta de control reúne los lugares donde el software se construye, se prueba y se ejecuta, y las herramientas que lo mueven entre ellos. Comprende cinco ambientes: Desarrollo, QA, Preproducción, Producción y Recuperación ante desastres. {\color{lafroxGris}Fuente: (Bases Administrativas, arts. 17.2–17.6 y Formulario T-14, p. 63); (Bases Técnicas Transversales, secciones 17.2–17.6, pp. 32–34).}} \tabularnewline*
  3.1.1 & Ambientes de Desarrollo y QA en la nube (H3) & Ambientes de Desarrollo y QA operativos. & La cadena de integración despliega todos los módulos en ambos ambientes. & SRE & E1 · Hasta el mes 5 (H3). \\
  3.1.2 & Ambientes de Preproducción y Producción en la nube (H3) & Ambientes de Preproducción y Producción operativos. & La misma versión del software pasa de QA a Preproducción y a Producción sin volver a compilarse; toda diferencia de escala entre Preproducción y Producción está declarada. & SRE & E1 · Hasta el mes 5 (H3). \\
  3.1.3 & Ambiente de recuperación ante desastres en la región secundaria (H3) & Ambiente de recuperación habilitado, con su procedimiento de conmutación. & El ambiente está habilitado en el H3. La prueba real se hace en 6.4.1. & SRE & E1 · Hasta el mes 5 (H3). \\
  3.1.4 & Cadena de integración y entrega continua, e infraestructura como código & Cadena operativa y repositorio de infraestructura como código. & Un cambio pasa por la cadena hasta QA sin intervención manual; los ambientes se pueden recrear desde el repositorio. & SRE & E1 · Mes 4. \\
  3.1.5 & Observabilidad de los doce módulos y de los sitios, con tableros para el CLIENTE (H3) & Tableros y alertas operativos. & Cada módulo y cada sitio publica sus indicadores con alertas configuradas. & SRE & E1 · Hasta el mes 5 (H3: «con observabilidad operativa»). \\
  \midrule
  \multicolumn{6}{>{\columncolor{lafroxBg}\RaggedRight\arraybackslash}p{\dimexpr\textwidth-2\tabcolsep\relax}}{\textbf{Cuenta 3.2 — Servicios de nube de la plataforma.} Esta cuenta reúne la puesta en marcha de los servicios de nube del T-11, su configuración según el plan de seguridad y su conexión con la observabilidad para uso del software. {\color{lafroxGris}Fuente: (Bases Administrativas, arts. 17.2–17.6 y Formulario T-14, p. 63); (Bases Técnicas Transversales, secciones 17.2–17.6, pp. 32–34).}} \tabularnewline*
  3.2.1 & Servicio de nube para la aplicación y los portales & Servicio configurado. & El software se despliega en él y cumple el plan de seguridad. & SRE & E1 · Mes 4. \\
  3.2.2 & Servicio de nube de base de datos, caché y mensajería & Servicio configurado. & La base de datos replica a la región secundaria y las colas entregan los mensajes sin duplicarlos. & SRE & E1 · Mes 4. \\
  3.2.3 & Servicio de nube de ingesta de IoT, telemetría cruda y plataforma analítica & Servicio configurado. & Las lecturas de un sensor de prueba llegan al almacén analítico. & SRE & E1 · Mes 4. \\
  3.2.4 & Servicio de nube de respaldo inmutable y réplica entre regiones & Servicio configurado, con su política de retención. & Una restauración de prueba recupera los datos. & SRE & E1 · Mes 4. \\
  3.2.5 & Servicio de nube de seguridad y gobierno, y detección en endpoints & Servicio configurado. & Los controles del plan de seguridad están activos y los equipos físicos reportan a la consola. & SEG & E1 · Mes 4. \\
  3.2.6 & Servicio de gestión de dispositivos móviles & Servicio configurado. & Un terminal de prueba se enrola, recibe la aplicación y se borra a distancia. & SRE & E1 · Antes del enrolamiento de los terminales (8.6.1). \\
  \midrule
  \multicolumn{6}{>{\columncolor{lafroxBg}\RaggedRight\arraybackslash}p{\dimexpr\textwidth-2\tabcolsep\relax}}{\textbf{Cuenta 3.3 — Servicios comunes.} Esta cuenta reúne los servicios comunes de la plataforma. La identidad, el registro auditable y la operación sin conexión (3.3.1–3.3.3) habilitan el comienzo de la construcción de los módulos de la Etapa 1. El motor de reglas, el repositorio documental y la generación mediante plantillas (3.3.4–3.3.6) se entregan en QA para la integración de la Etapa 1 (6.1.1). {\color{lafroxGris}Fuente: (Bases Administrativas, arts. 17.2–17.6 y Formulario T-14, p. 63); (Bases Técnicas Transversales, secciones 17.2–17.6, pp. 32–34).}} \tabularnewline*
  3.3.1 & Identidad y control de acceso, incluidos los conductores externos que rotan sin aviso & Identidad y control de acceso funcionando en QA. & Un conductor externo nuevo se habilita en el andén y pierde el acceso al terminar; la identidad funciona durante un corte de 24 horas. & SEG & E1 · Antes del H4. \\
  3.3.2 & Registro auditable de stock, lotes y cobros & Registro auditable en QA. & Toda operación de stock, lote y cobro queda registrada y se puede consultar. & ARQ & E1 · Antes del H4. \\
  3.3.3 & Operación sin conexión: 14 horas en terreno y 24 horas en el CD & Mecanismo de operación sin conexión en QA. & En Preproducción, una jornada de 14 horas sin señal y un corte de 24 horas del CD terminan sin pedidos perdidos ni duplicados. & ARQ & E1 · Antes del H4. \\
  3.3.4 & Parametrización y motor de reglas versionados & Consola de parámetros y motor de reglas en QA. & Un administrador cambia una regla sin recompilar; la versión queda con autor y fecha, y cada decisión registra la regla aplicada. & DES & E1 · Antes del H4. \\
  3.3.5 & Gestión documental con versionado, búsqueda y vista previa & Repositorio documental en QA. & Un documento se carga cifrado, conserva versiones y metadatos, respeta permisos, se encuentra por contenido o metadato, se previsualiza sin descarga y su integridad se verifica. & DES & E1 · Antes del H4. \\
  3.3.6 & Generación de documentos con plantillas administrables & Servicio de plantillas en QA. & Una plantilla editada por un administrador genera el documento en PDF, DOCX y XLSX con los datos de la operación. & DES & E1 · Antes del H4. \\
\end{tablalafrox}
